
\begin{remark}\label{fga-rem-2-10}\dcref{2.10}\source{fga-explained:page-0028}\uses{}
\begin{verbatim}
There isa dual versionof Yoneda's lemma, which willbe used
in ?3.2.1.Each object X of C definesa functor
                          Homc(X,-):C?^      (Set).
This can be viewed as the functor hx'- (C?p)?p -^ (Set);hence, from the usual
form of Yoneda's lemma applied to C?p for any two objectsX and Y we get a
canonicalbijectivecorrespondence between Homc(X, Y") and the set of natural
transformationsHomc(F, ?) -^ Homc(X, ?).

                             2.2. Group    objects
    In thissectionthe categoryC willhave finite
                                              products;we willdenote a terminal
objectby pt.
\end{verbatim}
\end{remark}


\begin{remark}\label{fga-rem-2-14}\dcref{2.14}\source{fga-explained:page-0030}\uses{}
\begin{verbatim}
Suppose thatC and V are categories with products and terminal
object pt^ and pt^^. Suppose that F: C ^ P is a functorthat preservesfinite
products,and G isa group object in C. The arrow e^: pt^ -^ G yieldsan arrow
Fee '?FptQ -^ FG; thiscan be composed with the inverseof the unique arrow
Fpt^ -^ pt^^,which is an isomorphism, because Fpt^ is a terminal object,to
get an arrow epc- P^v ~^ ^^- Analogously one uses Fm^: F{G x G) ^ FG
and the inverseof the isomorphism F{G x G) c^ FG x FG to definean arrow
niFG: FGxFG^       FG. Finallywe set ipc = Fie: FG -^ FG.
    We leaveitto the readerto check that thisgivesFG the structureof a group
object,and thisinducesa functorfrom the categoryof group objectson C to the
categoryof group objectson V.
     2.2.1. Actions of group objects. There isan obvious notionof leftaction
of a functorintogroups on a functorintosets.
\end{verbatim}
\end{remark}


\begin{remark}\label{fga-rem-2-23}\dcref{2.23}\source{fga-explained:page-0035}\uses{}
\begin{verbatim}
The terminology "C has discretegroup objects"isperhaps
        forC to have discretegroup objectswould be sufficient
misleading;                                                    to have discrete
objects,and that the functorA preservesfiniteproducts.
    However, fordiscretegroup objectsto be wellbehaved we need more than their
existence,we want Proposition2.22 to hold: and forthispurpose the conditions
of Definition2.20 seem to be optimal (exceptthat one does not need to assume
that C has allproducts;but thishypothesisissatisfied in allthe examples I have
in mind).

                  2.3. Sheaves in Grothendieck        topologies
    2.3.1. Grothendieck topologies. The readerisfamiliarwith the notion of
sheafon a topologicalspace. A presheafon a topologicalspace X can be considered
as a functor.Denote by Xd the categoryin which the objectsare the open subsets
of X, and the arrows are given by inclusions.Then a presheafof setson X is
a functorXci?^ -^ (Set);and thisisa sheaf when itsatisfies    appropriategluing
conditions.
    There aremore generalcircumstancesunder which we can ask whether a functor
isa sheaf.For example, considera functorF: (Top)?P -^ (Set);foreach topological
space X we can considerthe restriction  Fx to the subcategory Xd of (Top). We
say that F isa      on
               sheaf (Top) Fx if    isa sheafon X forallX.
    There is a very general notion of sheaf in a Grothendieck topology; in this
Sectionwe review thistheory.
    In a Grothendieck topology the "open sets"ofa space are maps intothisspace;
insteadof intersectionswe have to look at fiberedproducts,while unions play no
role.The axioms do  not describethe "open sets",but the coveringsof a space.
\end{verbatim}
\end{remark}


\begin{remark}\label{fga-rem-2-25}\dcref{2.25}\source{fga-explained:page-0036}\uses{}
\begin{verbatim}
In factwhat we have definedhereiswhat iscalleda pretopology
in [SGA4]; a pretopologydefinesa topology,and very different  pretopologiescan
definethe same topology.The point isthat the sheaf theory only depends on the
topology,and not on the pretopology.Two pretopologiesinducethe same topology
ifand only ifthey are equivalent,in the senseof Definition2.47.
    Despite itsunquestionabletechnicaladvantages,I do not findthe notion of
topology,as definedin [SGA4], very intuitive, so I preferto avoid itsuse (justa
questionof habit,undoubtedly).
    However, sieves^the objectsthat intervenein the definitionof a topology,are
quiteuseful,and willbe used extensively.
    Here are some examples of Grothendieck topologies.In what follows,a set
{Ui ^^ U} of functions,or morphisms of schemes, iscalledjointlysurjective
                                                                        when
the set-theoretic
                union of theirimages equals U.
\end{verbatim}
\end{remark}


\begin{lemma}\label{fga-lem-2-43}\dcref{2.43}\source{fga-explained:page-0042}\uses{}
\begin{verbatim}
IfF isseparated,the restriction
                                                function

                          HomE,F)    ?> Hom(h^^,F)
isinjective.
\end{verbatim}
\end{lemma}
\begin{proof}
\begin{verbatim}
Let us take two naturaltransformationscp.ip: S ^ F with the same
      in
image Hom(h^^,F),    an element T ^ U oi ST, and letus show          ^ U) =
                                                            that (/)(r
ip{T ^U)  e FT.
    Set U = {Ui -^ U}, and considerthe fiberedproducts T Xu Ui with their
projectionspi\ T XuUi ^T. SinceT Xf/[/^^ [/isin h^^(r Xf/ Ui) we have
           p*(/)(r^U) = cl){TXu Ui) = ^{T Xu Ui) = pl^iT ^ U).
Since {pii T Xu Ui -^ T} is a covering and F is a presheaf,we conclude that
    ^U) = V^(r^U) e ft, as desired.
(/)(r                                                                    D

     This concludesthe proof of Proposition2.42.                            D

    We concludewith a remark. Suppose that U = {Ui -^ U} and V = {Vj -^ U}
are coverings.Then U Xu V = {Ui Xu Vj ^ U} isslcovering.An arrow T ^ U
factorsthrough Ui XuVj ifand only ifitfactorsthrough Ui and through Uj. This
simple observationiseasilyseen to imply the followingfact.
\end{verbatim}
\end{proof}


\begin{proposition}\label{fga-pro-2-49}\dcref{2.49}\source{fga-explained:page-0043}\uses{fga-prop-2-42}
\begin{verbatim}
Let T and T' be two Grothendieck topologieson the same
categoryC. IfT issubordinateto T', then every sheafin T' isalsoa sheafin T.
    In particular,
                 two equivalenttopologieshave the same sheaves.
    The proof isimmediate from Propositions2.42and 2.48.
    In Grothendieck'slanguage what we have definedwould be calleda pretopology,
and two equivalentpretopologiesdefinethe same topology.
\end{verbatim}
\end{proposition}
\begin{proof}
\notready
\end{proof}


\begin{proposition}\label{fga-pro-2-54}\dcref{2.54}\source{fga-explained:page-0044}\uses{}
\begin{verbatim}
A representablefunctor (Top)?P -^ (Set)isa sheafin the
globalclassical
              topology.
    This amounts to sayingthat,given two topologicalspaces U and X, an open
covering{Ui C U), and continuousfunctions/^: Ui ^ X, with the property that
the restriction
              of fi and fj to Ui flUj coincideforalli and j, thereexistsa unique
continuousfunctionU ^ X whose restriction      Ui ^ X is fi. This is essentially
obvious (itboilsdown to the factthat,fora function,the property of being
        is localon the domain). For similarreasons,it is easy to show that a
continuous
representablefunctoron the category (Sch/5) over a base scheme 5 isa sheafin
the Zariskitopology.
    On the other hand the followingisnot easy at all:a scheme isa topological
space,togetherwith a sheafof ringsin the Zariskitopology.A priori,theredoes
not seem to be a reason why we should be ableto gluemorphisms of schemes in a
finertopology than the Zariskitopology.
\end{verbatim}
\end{proposition}
\begin{proof}
\notready
\end{proof}


\begin{remark}\label{fga-rem-2-56}\dcref{2.56}\source{fga-explained:page-0045}\uses{}
\begin{verbatim}
As we have already observed at the beginning of ?2.3.2,there
isa "wild"flattopology in which the coveringsare jointlysurjectivesets{Ui -^ U}
offlatmorphisms. However, thistopology isvery badly behaved; in particular,  not
allrepresentablefunctorsare sheaves.
    Take an integralsmooth curveU overan algebraically  closedfield,with quotient
fieldK and letVp = SpecOu^p forallclosedpointsp G U{k), as in Remark 1.14.
Then {V^ ^ [/} isa coveringin thiswild flattopology.
    Each Vp containsthe closedpointp, and Vp \ {p} = Spec K isthe genericpoint
of Vp-,furthermore Vp Xjj Vq = Vp iip = q, otherwiseVp Xjj Vq = Specif.
    We can form a (verynon-separated)scheme X by gluingtogether allthe Vp
along Spec K; then the embeddings Vp ^^ X and Vq ^^ X agree when restricted   to
VpXuVq,   so the givean element of Up ^x Vp whose two       in
                                                      images Up q^x{VpXuVq)
agree.However, thereisno morphism U r^ X whose restriction      to each Vp isthe
natural morphism Vp ^^ U. In fact,such a morphism would have to send each
closedpointp E U intop G V^ C X, and the genericpoint to the genericpoint;
but the resultingset-theoretic functionU ^ X isnot continuous,sinceallsubsets
of X formed by closedpointsare closed,while only the finitesetsare closedin U.
\end{verbatim}
\end{remark}


\begin{lemma}\label{fga-lem-2-61}\dcref{2.61}\source{fga-explained:page-0047}\uses{}
\begin{verbatim}
The sequence
                         0       AJ-.B^^^^Bi           B
isexact.
\end{verbatim}
\end{lemma}
\begin{proof}
\begin{verbatim}
The injectivity of / isclear,because B isfaithfully
                                                              flatover A. Also,
itisclearthat the image of / iscontainedin the kernelof ei ? 62,so we have only
to show that the kernelof ei ?62 iscontainedin the image of /.
     Assume that thereexistsa homomorphism of ^-algebrasg: B ^ A {inother
words, assume that the morphism V ^^ U has a section).Then the composite
g o f: ^ ^ ^ isthe identity.Take an element b G ker(ei? 62);by definition, this
means that b^l = l^bin B(g)^B. By applyingthe homomorphism g(g)id^: B(g^
38                       2.CONTRAVARIANT FUNCTORS

B -^ A^A B = B to both members ofthe equalitywe obtainthat f{gb) = b,hence
be im/.
    In general,therewillbe no sectionU ^ V; however,suppose thatthereexistsa
         flatA algebraA^^ A ^such thatthe homomorphism j^idiA'?A! -^ B<S>A
faithfully
obtained by base change has a sectionB <S>A' ^ A' 8iSbefore.Set B' = B <S>A'.
Then thereisa naturalisomorphism of ^'-algebrasB' (g)^/ B' :^ {B <S>aB) ?a A',
making the diagram
                         /(8)id^
                              ^B'                           ->B' ?A' B'


                                       (ei-e2
                                ^B'                 ->{B ?A B) ?A A'

commutative. The top row is exact,because of the existenceof a section,and so
the bottom row isexact.The thesisfollows,because A' isfaithfullyflatover A.
    But to findsuch homomorphism A ^ A' \tisenough to setA' = B\ the product
B'^A A' -^ A' definedhy b^b' ^-^bb'givesthe desiredsection.In geometricterms,
the diagonalV ^V XuV givesa sectionof the firstprojectionV XuV ^V.          D

    To finishthe proof of Theorem 2.55 in the case that X is affine,recallthat
morphisms of schemes U ^ X, V ^^ X and F Xf/ F ^ X correspond to ring
homomorphisms R ^^ A, R ^ B and R ^ B (g)^B; then the resultisimmediate
from the lemma above. This proves that hx isa sheafwhen X isaffine.
    IfX isnot necessarilyaffine, write X = DiXi as a union of affineopen sub-
schemes.
    Let us show that hx isseparated.Given a coveringV ^ U, taketwo morphisms
f,g: U ^^ X such that the two composites V ^^ U ^^ X are equal.SinceV ^ U is
surjective,                                 so we can set Ui f-^X, = g-^X,,
           / and g coincideset-theoretically,
and callVi the inverseimages of Ui in V. The two composites

                                              flu.
                            V,        u.         i X,
                                              9\u,

coincide,and Xi isaffine;hence / \u,= 9 \u,forallz,so f = g, as desired.
    To complete the proof,suppose that g: F ^ X isa morphism with the property
that the two composites
                                        Pi-i
                          VXrjV                 V       X
                                        Pr2
are equal;we need to show that g factorsthrough U. The morphism V ^ U is
surjective,so,from Lemma 2.62 below,g factorsthrough U set-theoretically.
                                                                      Since
U has the quotienttopology induced by the morphism V ^ U (Proposition1.13),
we get that the resultingfunction/: U ^ X iscontinuous.
    Set Ui = f~^Xi and Vi = g~^Vi foralli.The composites
                                  Pi-i        q\v
                      Vi XrrV              Vi ^ Vi           Xi
                                  Pr2
2.3.SHEAVES IN GROTHENDIECK TOPOLOGIES                     39

coincide,and Xi is affine,so ^ |v^:V^ ^ X factorsuniquelythrough a morphism
fi'.Ui^^Xi.  We have
                                         - Ui D Uj ?> X,
                      fi \u^nUj= fj \u^r]UJ
because hx isseparated;hence the figluetogetherto givethe desiredfactorization
\end{verbatim}
\end{proof}
